\documentclass[11pt]{article}
\usepackage{amsmath,amssymb,amsthm,geometry,hyperref}
\geometry{margin=25mm}
\newtheorem{theorem}{Theorem}
\newcommand{\Tr}{\operatorname{Tr}}
\begin{document}
\title{A cubic auxiliary-spin bound on third-order anticoherent subspaces}
\author{Exact catalogue calculation}
\date{3 November 2026}
\maketitle
\begin{theorem}
For every integer $N\ge19$ and AC order $t\ge3$,
\[
\boxed{6K(N,t)\le N-4.}
\]
In particular, $K(19,3)\le2$ and $K(20,3)\le2$.
\end{theorem}
\begin{proof}
Let $j=N/2$, and let $P$ project onto a $k$-dimensional third-order AC
subspace of $V_j$. Polarization gives $PT_{\ell q}P=0$ for every tensor
of rank $1\le\ell\le3$. Every ordered spin polynomial $A$ of degree at
most three therefore satisfies $PAP=(\Tr A)/(2j+1)P$.

Introduce an auxiliary spin $5/2$, and set
\[
X=\sum_{a=x,y,z}J_a\otimes S_a,
\quad F_j(x)=-(2x+3j+8)(2x-3j+5)(2x-5j).
\]
For $p=0,\ldots,5$, the total-spin sectors $L=j-5/2+p$ have dimensions
$2j-4+2p$ and the following coupling eigenvalues and polynomial values:
\[
\begin{array}{c|c|c}
p&x_L&F_j(x_L)\\\hline
0&-5(j+1)/2&40j(2j-3)(2j+1)\\
1&-3j/2-4&0\\
2&-j/2-9/2&-12(j+1)(2j-1)(2j+3)\\
3&j/2-4&-16j(j+2)(2j+3)\\
4&3j/2-5/2&0\\
5&5j/2&0
\end{array}
\]
Here $x_L=[L(L+1)-j(j+1)-35/4]/2$. The weighted sum gives
\[
\frac{\Tr F_j(X)}{6(2j+1)}
=c_j=\frac{20}{3}j(j-9)(2j-1)>0\qquad(j\ge19/2).
\]
Set $W=F_j(X)-c_jI$. Since $F_j(X)$ is invariant under joint rotations,
its partial trace over $V_j$ is scalar on the auxiliary irreducible
space. The above trace fixes that scalar and yields
\[
\Tr_{V_j}W=0.
\]
Every sector except $L=j-5/2$ has a strictly negative eigenvalue of $W$.
The remaining eigenvalue is positive, either from $\Tr W=0$ or directly:
\[
F_j(-5(j+1)/2)-c_j
=\frac{20}{3}j(j+1)(22j-27)>0.
\]
Thus $W$ is nonsingular with positive index $N-4$ and negative index
$5(N+2)$.

The degree-three compression identity gives
\[
(P\otimes I_6)W(P\otimes I_6)
=\frac{P}{2j+1}\otimes\Tr_{V_j}W=0.
\]
Hence $\operatorname{Ran}P\otimes V_{5/2}$ is a $6k$-dimensional totally
isotropic subspace. Projection into the positive spectral space is
injective: a nonzero vector with zero positive component would have a
strictly negative quadratic form. Therefore $6k\le N-4$.
Every higher-order AC subspace is third-order AC as well.
\end{proof}
At $N=19$, this gives $6k\le15$; at $N=20$, it gives $6k\le16$.
The upper bound excludes dimension three at both sizes, without assuming
real amplitudes, a Dicke support pattern or constellation symmetry.

\paragraph{Exact verification.}
The accompanying verifier reconstructs the physical magnetic blocks of
$J\cdot S$ at $N=19,20$, checks their characteristic polynomials and
Hermiticity with the exact physical metric, and verifies every partial-
trace entry of $W$ with rational arithmetic. The general weighted-trace
identity is also checked at five exact points after its degree-four
bound is established. Numerical searches play no role in the proof.
\end{document}
